% Part: incompleteness
% Chapter: arithmetization-syntax
% Section: coding-symbols

\documentclass[../../../include/open-logic-section]{subfiles}

\begin{document}

\olfileid[id]{inc}{art}{cod}
\olsection{Pengodean Simbol}

Bahasa dasar~$\Lang L$ bagi logika orde pertama menggunakan simbol-simbol
\[
\lfalse \quad \lnot \quad \lor \quad \land \quad \lif \quad \lforall
\quad \lexists \quad \eq \quad ( \quad ) \quad ,
\]
bersama himpunan variabel dan !!{constant}s yang !!{enumerable}, serta
himpunan !!{function}s dan !!{predicate}s beraritas sembarang yang
!!{enumerable}. Kita dapat memasangkan \emph{kode}
dengan setiap simbol ini
sedemikian rupa sehingga setiap simbol memperoleh sebuah bilangan unik sebagai
kodenya dan tidak ada dua simbol berbeda yang memperoleh bilangan yang sama.
Kita tahu bahwa hal ini mungkin karena himpunan semua simbol !!{enumerable},
sehingga terdapat !!a{bijection} antara himpunan itu dan himpunan bilangan asli.
Namun, kita ingin memastikan bahwa simbolnya (serta beberapa informasi tentang
simbol itu, misalnya aritas !!a{function}) dapat ditemukan kembali dari kodenya
dengan cara yang dapat dihitung. Tentu saja ada banyak cara untuk melakukannya.
Berikut salah satunya, yang menggunakan fungsi-fungsi rekursif primitif.
(Ingat bahwa $\tuple{n_0, \dots, n_k}$ adalah bilangan yang mengodekan barisan
bilangan $n_0$, \dots, $n_k$.)

\begin{defn}
Jika $s$ adalah simbol~$\Lang L$, definisikan \emph{kode simbol}~$\scode s$
sebagai berikut:
\begin{enumerate}
\item Jika $s$ termasuk simbol logis, $\scode s$ diberikan oleh tabel berikut:
\[
\begin{array}{cccccccccc}
  \lfalse & \lnot & \lor & \land & \lif & \lforall \\
\tuple{0, 0} & \tuple{0, 1} & \tuple{0, 2} & \tuple{0, 3} &
\tuple{0, 4} & \tuple{0, 5} \\
\lexists & \eq & ( & ) & ,\\
\tuple{0, 6} & \tuple{0, 7} &
\tuple{0, 8} & \tuple{0, 9} & \tuple{0, 10}
\end{array}
\]
\item Jika $s$ adalah variabel ke-$i$, $\Obj v_i$, maka
  $\scode s = \tuple{1, i}$.
\item Jika $s$ adalah !!{constant} ke-$i$,~$\Obj c_i$, maka
  $\scode s = \tuple{2, i}$.
\item Jika $s$ adalah !!{function} ke-$i$ yang beraritas $n$,~$\Obj f_i^n$,
  maka $\scode s = \tuple{3, n, i}$.
\item Jika $s$ adalah !!{predicate} ke-$i$ yang beraritas $n$,~$\Obj P_i^n$,
  maka $\scode s = \tuple{4, n, i}$.
\end{enumerate}
\end{defn}

\begin{prop}
Relasi-relasi berikut rekursif primitif:
\begin{enumerate}
\item $\fn{Fn}(x, n)$ jika dan hanya jika $x$ adalah kode $\Obj f^n_i$ untuk
  suatu~$i$, yakni $x$ adalah kode suatu !!{function} beraritas $n$.
\item $\fn{Pred}(x, n)$ jika dan hanya jika $x$ adalah kode $\Obj P^n_i$ untuk
  suatu~$i$, atau $x$ adalah kode $\eq$ dan $n = 2$, yakni $x$ adalah kode
  suatu !!{predicate} beraritas $n$.
\end{enumerate}
\end{prop}

\begin{defn}
Jika $s_0, \dots, s_{n-1}$ adalah barisan simbol, \emph{bilangan
  G\"odel}-nya ialah $\tuple{\scode{s_0}, \dots, \scode{s_{n-1}}}$.
\end{defn}

\begin{explain}
Perhatikan bahwa \emph{kode} dan \emph{bilangan G\"odel} merupakan dua hal
yang berbeda. Sebagai contoh, variabel~$\Obj v_5$ memiliki kode~$\scode{\Obj
  v_5} = \tuple{1, 5} = 2^2\cdot 3^6$. Namun, variabel~$\Obj v_5$ yang
ditinjau sebagai suku juga merupakan barisan simbol (dengan panjang~$1$).
\emph{Bilangan G\"odel}~$\Gn{\Obj v_5}$ dari \emph{suku}~$\Obj v_5$ ialah
$\tuple{\scode{\Obj v_5}} = 2^{\scode{\Obj v_5} + 1} =
2^{2^2\cdot 3^6 + 1}$.
\end{explain}

\begin{ex}
Ingat bahwa jika $k_0$, \dots, $k_{n-1}$ adalah barisan bilangan, maka kode
barisan $\tuple{k_0, \dots, k_{n-1}}$ dalam pengodean dengan pangkat bilangan
prima ialah
\[
2^{k_0+1}\cdot3^{k_1+1}\cdot \dots \cdot p_{n-1}^{k_{n-1}+1},
\]
dengan $p_i$ bilangan prima ke-$i$ (dimulai dengan $p_0 = 2$). Jadi, sebagai
contoh, formula $\eq[\Obj v_0][\Obj 0]$, atau, jika dituliskan lebih lengkap,
${\eq}(\Obj v_0,\Obj c_0)$, memiliki bilangan G\"odel
\[
\tuple{\scode{\eq},\scode{(},\scode{\Obj v_0},\scode{,}, \scode{\Obj
    c_0},\scode{)}}.
\]
Di sini, $\scode{\eq}$ ialah $\tuple{0,7} = 2^{0+1}\cdot
3^{7+1}$, $\scode{\Obj v_0}$ ialah $\tuple{1,0} = 2^{1+1}\cdot3^{0+1}$,
dan seterusnya. Jadi, $\Gn{=(\Obj v_0,\Obj c_0)}$ ialah
\begin{multline*}
2^{\scode{=} + 1}\cdot 3^{\scode{(}+1}\cdot 5^{\scode{\Obj v_0}+1}
\cdot 7^{\scode{,} + 1} \cdot 11^{\scode{\Obj c_0}+1} \cdot
13^{\scode{)}+1} = \\
2^{2^1\cdot 3^8 + 1}\cdot 3^{2^1\cdot 3^9+1}\cdot 5^{2^2\cdot 3^1+1}
\cdot 7^{2^1\cdot 3^{11} + 1} \cdot 11^{2^3\cdot3^1+1} \cdot
13^{2^1\cdot3^{10}+1} = \\
2^{13\,123}\cdot 3^{39\,367}\cdot 5^{13}\cdot 7^{354\,295}\cdot11^{25}\cdot13^{118\,099}.
\end{multline*}
\end{ex}

\end{document}
